\section{Locally free sheaves}
\label{section-locally-free}

\noindent
Let $(X, \mathcal{O}_X)$ be a ringed space.
Our conventions allow (some of) the stalks $\mathcal{O}_{X, x}$
to be the zero ring. This means we have to be a little careful
when defining the rank of a locally free sheaf.

\begin{definition}
\label{definition-locally-free}
Let $(X, \mathcal{O}_X)$ be a ringed space.
Let $\mathcal{F}$ be a sheaf of $\mathcal{O}_X$-modules.
\begin{enumerate}
\item We say $\mathcal{F}$ is {\it locally free} if for every
point $x \in X$ there exist a set $I$ and an open
neighbourhood $x \in U \subset X$
such that $\mathcal{F}|_U$ is isomorphic to
$\bigoplus_{i \in I} \mathcal{O}_X|_U$ as an $\mathcal{O}_X|_U$-module.
\item We say $\mathcal{F}$ is {\it finite locally free} if we may
choose the index sets $I$ to be finite.
\item We say $\mathcal{F}$ is {\it finite locally free of rank $r$}
if we may choose the index sets $I$ to have cardinality $r$.
\end{enumerate}
\end{definition}

\noindent
A finite direct sum of (finite) locally free sheaves is (finite)
locally free. However, it may not be the case that an infinite direct
sum of locally free sheaves is locally free.

\begin{lemma}
\label{lemma-locally-free-quasi-coherent}
Let $(X, \mathcal{O}_X)$ be a ringed space.
Let $\mathcal{F}$ be a sheaf of $\mathcal{O}_X$-modules.
If $\mathcal{F}$ is locally free then it is quasi-coherent.
\end{lemma}

\begin{proof}
Omitted.
\end{proof}

\begin{lemma}
\label{lemma-pullback-locally-free}
Let $f : (X, \mathcal{O}_X) \to (Y, \mathcal{O}_Y)$
be a morphism of ringed spaces. If $\mathcal{G}$ is
a locally free $\mathcal{O}_Y$-module, then
$f^*\mathcal{G}$ is a locally free $\mathcal{O}_X$-module.
\end{lemma}

\begin{proof}
Omitted.
\end{proof}

\begin{lemma}
\label{lemma-rank}
Let $(X, \mathcal{O}_X)$ be a ringed space.
Suppose that the support of $\mathcal{O}_X$ is $X$,
i.e., all stalks of $\mathcal{O}_X$ are nonzero rings.
Let $\mathcal{F}$ be a locally free sheaf of $\mathcal{O}_X$-modules.
There exists a locally constant function
$$
\text{rank}_\mathcal{F} :
X \longrightarrow \{0, 1, 2, \ldots\}\cup\{\infty\}
$$
such that for any point $x \in X$ the cardinality of any
set $I$ such that $\mathcal{F}$ is isomorphic to
$\bigoplus_{i\in I} \mathcal{O}_X$ in a neighbourhood
of $x$ is $\text{rank}_\mathcal{F}(x)$.
\end{lemma}

\begin{proof}
Under the assumption of the lemma the cardinality of $I$ can be read off
from the rank of the free module $\mathcal{F}_x$ over the nonzero ring
$\mathcal{O}_{X, x}$, and it is constant in a neighbourhood of $x$.
\end{proof}

\begin{lemma}
\label{lemma-map-finite-locally-free}
Let $(X, \mathcal{O}_X)$ be a ringed space. Let $r \geq 0$.
Let $\varphi : \mathcal{F} \to \mathcal{G}$ be a map of finite
locally free $\mathcal{O}_X$-modules of rank $r$.
Then $\varphi$ is an isomorphism if and only if $\varphi$
is surjective.
\end{lemma}

\begin{proof}
Assume $\varphi$ is surjective. Pick $x \in X$.
There exists an open neighbourhood $U$ of $x$
such that both $\mathcal{F}|_U$ and $\mathcal{G}|_U$ are
isomorphic to $\mathcal{O}_U^{\oplus r}$.
Pick lifts of the free generators of $\mathcal{G}|_U$
to obtain a map $\psi : \mathcal{G}|_U \to \mathcal{F}|_U$
such that $\varphi|_U \circ \psi = \text{id}$.
Hence we conclude that the map
$\Gamma(U, \mathcal{F}) \to \Gamma(U, \mathcal{G})$
induced by $\varphi$ is surjective. Since both
$\Gamma(U, \mathcal{F})$ and $\Gamma(U, \mathcal{G})$
are isomorphic to $\Gamma(U, \mathcal{O}_U)^{\oplus r}$
as an $\Gamma(U, \mathcal{O}_U)$-module we may
apply Algebra, Lemma \ref{algebra-lemma-fun} to see that
$\Gamma(U, \mathcal{F}) \to \Gamma(U, \mathcal{G})$
is injective. This finishes the proof.
\end{proof}

\begin{lemma}
\label{lemma-direct-summand-of-locally-free-is-locally-free}
Let $(X, \mathcal{O}_X)$ be a ringed space. If all stalks $\mathcal{O}_{X, x}$
are local rings, then any direct summand of a finite locally free
$\mathcal{O}_X$-module is finite locally free.
\end{lemma}

\begin{proof}
Assume $\mathcal{F}$ is a direct summand of the finite locally free
$\mathcal{O}_X$-module $\mathcal{H}$.
Let $x \in X$ be a point. Then $\mathcal{H}_x$ is a finite free
$\mathcal{O}_{X, x}$-module.
Because $\mathcal{O}_{X, x}$ is local, we see that
$\mathcal{F}_x \cong \mathcal{O}_{X, x}^{\oplus r}$ for some $r$, see
Algebra, Lemma \ref{algebra-lemma-finite-projective}.
By Lemma \ref{lemma-finite-presentation-stalk-free}
we see that $\mathcal{F}$ is free of rank $r$ in an open neighbourhood of $x$.
(Note that $\mathcal{F}$ is of finite presentation as a summand of
$\mathcal{H}$.)
\end{proof}





